\documentclass[10pt,a4paper]{article}
\usepackage[margin=2.25cm]{geometry}
\usepackage[T1]{fontenc}
\usepackage{lmodern}
\usepackage{amsmath,amssymb,amsthm,bm,mathtools}
\usepackage{graphicx,booktabs,array,tabularx,longtable}
\usepackage{xcolor}
\usepackage[numbers,sort&compress]{natbib}
\usepackage[colorlinks=true,allcolors=blue!45!black]{hyperref}
\usepackage{microtype}
\usepackage{placeins}
\graphicspath{{figures/}}
\newcommand{\F}{\bm F}
\newcommand{\C}{\bm C}
\newcommand{\E}{\bm E}
\newcommand{\Sg}{\bm S}
\newcommand{\I}{\bm I}
\newcommand{\ee}{\bm e}
\newcommand{\ssv}{\bm s}
\newcommand{\qq}{\bm q}
\newcommand{\dd}{\mathrm d}
\newcommand{\R}{\mathbb R}
\newcommand{\GL}{\mathrm{GL}^{+}(2)}
\newcommand{\cof}{\operatorname{cof}}
\newcommand{\tr}{\operatorname{tr}}
\newcommand{\diag}{\operatorname{diag}}
\newcommand{\doi}[1]{\href{https://doi.org/#1}{doi:#1}.}
\newtheorem{proposition}{Proposition}
\newtheorem{remark}{Remark}
\setlength{\emergencystretch}{2em}
\newcolumntype{Y}{>{\raggedright\arraybackslash}X}
\title{Learned polyconvex constitutive models and hyperreduced micromechanics\protect\\
for anisotropic finite-strain homogenization}
\author{Working manuscript --- author list and affiliations to be confirmed}
\date{6 September 2026\\\small Evidence-based draft v0.1; full bibliography audit in progress}

\begin{document}
\maketitle
\begin{abstract}
Replacing microscopic boundary-value problems by inexpensive constitutive maps can substantially reduce the cost of finite-strain multiscale analysis. However, interpolation accuracy, mechanical admissibility, and computational efficiency are different requirements. This work investigates their interaction for an anisotropic porous hyperelastic material using two families of approximations: projection-based hyperreduced micromechanics and learned macroscopic constitutive laws. For the latter, we formulate paired directional features with trainable powers and determinant exponents, coupled to convex nondecreasing neural or spline networks. The construction preserves objectivity and two-dimensional polyconvexity, while affine determinant normalization separates the stress-free reference condition from volumetric stiffness. A separate sufficient condition addresses global nonnegative energy for the selected trained parameters. A periodic cell with a rotated elliptical void first tests constitutive accuracy and mechanical consistency. The selected ICNN and spline-ICKAN models attain aggregate held-out stress errors of approximately $0.65\%$. A closed strain cycle exposes nonconservative work in direct stress regression, whereas a search for negative rank-one curvature in an unconstrained energy network is reported with its reference-validation limitations. A tensile coupon then evaluates structural errors and online cost against a nested full-order reference. The polyconvex models produce displacement errors below $0.43\%$, while optimized hyperreduced micromechanics is more accurate but generally more expensive. All compared models converge on this nominal load path. The results support accurate, inexpensive polyconvex surrogates for this material and deformation domain, without establishing unrestricted extrapolation or universality across microstructures.
\end{abstract}
\noindent\textbf{Keywords:} computational homogenization; hyperelasticity; polyconvex neural networks; input-convex splines; empirical cubature; nonlinear model reduction.

\section{Introduction}
\label{sec:intro}
In first-order computational homogenization, a macroscopic constitutive evaluation is obtained from the solution of a microscopic boundary-value problem. This avoids prescribing an effective constitutive equation, but transfers a considerable computational burden to each macroscopic integration point. For nonlinear problems, the microscopic response and its tangent may be needed repeatedly during every macroscopic load increment. Reducing this burden without obscuring the mechanical content of the resulting model is the central concern of this work.

Two complementary strategies are considered. Projection-based methods compress the microscopic displacement field and continue to impose reduced equilibrium. The reduced model multiscale method of Yvonnet and He~\cite{yvonnet2007} is an early finite-strain example of this approach. For nonlinear finite element models, reducing the number of unknowns does not remove the cost of evaluating their residuals. Empirical cubature addresses this second cost by selecting integration contributions and associated weights~\cite{hernandez2017}. Nonlinear approximation manifolds can further reduce the number of online coordinates; the PROM--ANN construction of Barnett et al.~\cite{barnett2023} learns a closure in a reduced latent space, and manifold-adaptive empirical cubature has recently been developed for nonlinear finite element models~\cite{hernandez2026}. These ingredients serve here as established building blocks, not as newly introduced reduction algorithms.

The alternative is to replace microscopic equilibrium by a learned macroscopic constitutive relation. As'ad et al.~\cite{asad2022} explain why direct stress regression need not define a conservative material, even when its training data are hyperelastic, and construct energy-based models using objective strains. Their convexity condition is imposed in Green--Lagrange strain; the associated stability discussion explicitly distinguishes tensile and compressive states. It must not be relabelled as a general polyconvexity certificate. The FEANN framework of Kalina et al.~\cite{kalina2023} combines physics-constrained energy networks with automated multiscale data mining. These works also motivate distinguishing the selection of informative deformation states from the approximation of their constitutive response.

Polyconvexity supplies a different architectural constraint. In the framework introduced by Ball~\cite{ball1977}, the energy is represented by a convex function of deformation minors. Klein et al.~\cite{klein2022} implement polyconvex anisotropic neural energies using invariant-based and deformation-gradient-based constructions. The latter uses data augmentation for approximate objectivity; the former is objective by construction. Linden et al.~\cite{linden2023} give compatible reference-state and growth corrections, while treating global nonnegative energy as a separate question. More recently, Klein et al.~\cite{klein2026} develop objective polyconvex models based on triclinic invariants, with optional group symmetrization. Consequently, the absence of an imposed material symmetry group is not by itself a novelty claim for the present formulation. Input-convex Kolmogorov--Arnold networks provide another choice of convex nondecreasing core: Thakolkaran et al.~\cite{thakolkaran2025} develop such models for isotropic compressible hyperelasticity, including symbolic regression and finite element applications.

The unresolved practical issue addressed here is narrower than universal constitutive approximation. A mathematically admissible feature set can be too restrictive for a particular anisotropic homogenized response. Increasing network width cannot, in general, remove a restriction introduced before the network. Conversely, removing convexity constraints may improve a fit while removing a constitutive guarantee. We therefore study a feature family that couples directional distortions to the determinant, and compare its predictive and mechanical behavior with unrestricted learned models and with reduced micromechanics.

The contributions of this study are threefold. First, a paired-direction construction provides trainable polyconvex features whose reference derivatives remain isotropic, allowing a single affine determinant correction to normalize stress. Second, the same construction is implemented with a smooth ICNN and a $C^2$ convex spline core, with analytical properties separated from audits of particular trained parameters. Third, constitutive and structural tests are organized to answer different questions: the cell assesses approximation and mechanical consistency, whereas the coupon measures error propagation, solver behavior, and online cost against the same nested reference. The present evidence concerns one microstructure; the two tests are not presented as independent demonstrations across materials.

\section{Finite-strain homogenization and the reference problem}
\label{sec:fom}
\subsection{Kinematics and work-conjugate variables}
We consider an in-plane deformation gradient $\F\in\GL$, with
\begin{equation}
 \C=\F^T\F,\qquad J=\det\F>0,\qquad
 \E=\tfrac12(\C-\I).
\end{equation}
Plane strain means that the three-dimensional gradient is $\diag(\F,1)$; it does not mean that the out-of-plane stress vanishes. The engineering strain vector and its work-conjugate stress vector are
\begin{equation}
 \ee=(E_{11},E_{22},2E_{12})^T,\qquad
 \ssv=(S_{11},S_{22},S_{12})^T,\qquad
 \Sg:\dd\E=\ssv\cdot\dd\ee.
 \label{eq:voigt}
\end{equation}
For an objective differentiable energy $W$, the second Piola stress and constitutive tangent are $\ssv=\partial_{\ee}W$ and $\bm D=\partial^2_{\ee\ee}W$. The factors in Eq.~\eqref{eq:voigt} are essential for both energy-cycle tests and derivative checks.

\subsection{Periodic microscopic equilibrium}
Let $Y$ denote a reference periodic cell and $Y_s\subset Y$ its solid part. A homogeneous macroscopic gradient $\overline{\F}$ induces the microscopic deformation
\begin{equation}
 \bm\varphi(\bm X)=\overline{\F}\bm X+\bm w(\bm X),\qquad
 \F_\mu=\overline{\F}+\nabla\bm w,
\end{equation}
where $\bm w$ is periodic on opposite faces. A constraint removes its translational null space, and the void surface is traction-free. On a differentiable equilibrium branch, the discrete fluctuation coordinates $\bm a^*(\overline{\F})$ satisfy
\begin{equation}
 \bm r(\bm a,\overline{\F})=\partial_{\bm a}\Pi_\mu(\bm a,\overline{\F})=\bm0,
 \qquad
 \Pi_\mu=\int_{Y_s}\psi_\mu(\F_\mu)\,\dd V.
 \label{eq:microeq}
\end{equation}
The branch energy and homogenized first Piola stress are
\begin{equation}
 \overline W=\frac{\Pi_\mu(\bm a^*,\overline{\F})}{|Y|},\qquad
 \overline{\bm P}=\frac{1}{|Y|}\int_{Y_s}\bm P_\mu\,\dd V,
 \qquad \overline{\Sg}=\overline{\F}^{-1}\overline{\bm P}.
 \label{eq:homstress}
\end{equation}
The denominator includes the void. Averaging microscopic second Piola stress instead of first Piola stress would generally give a different quantity and is not used. Stationarity in Eq.~\eqref{eq:microeq} eliminates the fluctuation derivative in the first derivative of $\overline W$. Periodic kinematics and matching boundary tractions give the corresponding macroscopic work identity.

The implementation queries the constitutive law through $\overline{\F}=\sqrt{\I+2\E}$, exploiting objectivity to remove a superposed macroscopic rotation. The microscopic Newton method follows an equilibrium branch; convergence does not establish global minimization of $\Pi_\mu$. In particular, no general assertion that homogenization preserves polyconvexity is required here.

\subsection{Consistent tangent and the scope of the reference}
If $\bm r_{,\bm a}$ is nonsingular on the selected branch, implicit differentiation gives
\begin{equation}
 \frac{\dd\bm a^*}{\dd\ee}=-\bm r_{,\bm a}^{-1}\bm r_{,\ee},\qquad
 \bm D=\ssv_{,\ee}-\ssv_{,\bm a}\bm r_{,\bm a}^{-1}\bm r_{,\ee}.
 \label{eq:ift}
\end{equation}
Every derivative in this expression must refer to the stress and residual actually evaluated. The reference and reduced implementations are checked against perturbations that re-solve microscopic equilibrium. These checks validate linearization, not a change of homogenization assumptions.

The nested FOM--FE$^2$ solution is the reference for surrogate errors at fixed macro and micro discretizations. It is not a direct numerical simulation of all pores in the coupon. Homogenization error, boundary-layer effects, and mesh-discretization error therefore remain conceptually distinct from the surrogate errors reported below.

\section{Learned constitutive laws}
\label{sec:learned}
\subsection{What is changed between the four learned models?}
\label{sec:tiers}
Table~\ref{tab:tiers} compares the actual model families used in this study. It is deliberately separate from the literature comparison: a property of our implementation is not automatically attributed to an entire class of published methods.
\begin{table}[htbp]
\centering\small
\caption{Learned constitutive models in the present comparison. ``Not imposed'' is not synonymous with ``fails at every state''.}
\label{tab:tiers}
\begin{tabularx}{\linewidth}{lYYYY}
\toprule
 & Regression & Free & ICNN & ICKAN\\
\midrule
Network input & $E_{11},E_{22},\gamma_{12}$ & $C_{11}-1,C_{22}-1,C_{12},J-1$ & 32 paired features & 32 paired features\\
Network output & Three stresses & Scalar energy & Scalar energy core & Scalar energy core\\
Potential structure & Not imposed & Exact & Exact & Exact\\
Objectivity & Objective strain input and material stress output & Exact & Exact & Exact\\
Reference energy/stress & Not imposed & Exact normalization & Exact normalization & Exact normalization\\
Polyconvexity & Not imposed & Not imposed & By construction & By construction\\
Global $W\geq0$ & No potential & Not guaranteed & Separate saved-weight bound & Separate saved-weight bound\\
\bottomrule
\end{tabularx}
\end{table}

Regression directly approximates $\ee\mapsto\ssv$ with three outputs. The Free model instead defines an unconstrained energy network $\widetilde W(\ee)$ and applies
\begin{equation}
 W_{\mathrm{free}}(\ee)=\widetilde W(\ee)-\widetilde W(\bm0)
 -\partial_{\ee}\widetilde W(\bm0)\cdot\ee.
 \label{eq:free}
\end{equation}
This enforces zero reference energy and stress without enforcing polyconvexity. The four Free inputs are redundant coordinates of the same three-dimensional strain state; they are not four independent loading parameters. Its use of material-coordinate components of $\C$ does not impose isotropy. Reference-axis changes must transform the material directions together with the tensor components.

\subsection{Paired directional features}
For each feature, let $\bm d_i=(\cos\theta_i,\sin\theta_i)^T$ and $\bm d_i^\perp=(-\sin\theta_i,\cos\theta_i)^T$. Define
\begin{equation}
 z_i(\F,J)=\frac{\|\F\bm d_i\|^{2p_i}}{p_i J^{b_i}}
          +\frac{\|\F\bm d_i^\perp\|^{2q_i}}{q_i J^{c_i}},
 \label{eq:features}
\end{equation}
with
\begin{equation}
 p_i,q_i\geq\tfrac12,\qquad
 0\leq b_i\leq2p_i-1,\qquad 0\leq c_i\leq2q_i-1.
 \label{eq:featureconstraints}
\end{equation}
The parameters $\theta_i,p_i,q_i,b_i,c_i$ are learned subject to Eq.~\eqref{eq:featureconstraints}. The powers use a shifted softplus parametrization, and the determinant exponents use sigmoid fractions of their admissible upper bounds. The paired terms need not have equal exponents. Consequently, pairing perpendicular directions does not impose invariance under their interchange.

\begin{proposition}
Each feature in Eq.~\eqref{eq:features} is objective and convex in the independent variables $(\F,J)$ on $J>0$. Its value and engineering-strain gradient at the identity are
\begin{equation}
 z_i^0=\frac1{p_i}+\frac1{q_i},\qquad
 \partial_{\ee}z_i\big|_{\bm0}=\rho_i(1,1,0)^T,
 \qquad \rho_i=2-\frac{b_i}{p_i}-\frac{c_i}{q_i}.
 \label{eq:referencefeatures}
\end{equation}
\end{proposition}
\begin{proof}
The map $\|\bm x\|^a J^{-b}$ is jointly convex for $a\geq1$ and $0\leq b\leq a-1$, as shown in Appendix~\ref{app:convex}. Composition with the linear maps $\F\mapsto\F\bm d_i$ and $\F\mapsto\F\bm d_i^\perp$ proves convexity. Norms are unchanged by left rotations, proving objectivity. At the identity, differentiating the two numerators yields $2\bm d_i\otimes\bm d_i$ and $2\bm d_i^\perp\otimes\bm d_i^\perp$ with respect to $\E$. Their sum is $2\I$. The determinant factors contribute $-(b_i/p_i+c_i/q_i)\I$, giving Eq.~\eqref{eq:referencefeatures}.
\end{proof}

Positive feature scaling and subtraction of a reference feature vector improve numerical conditioning without changing these properties. They are constants with respect to deformation for every frozen parameter set. Learning the directions does not amount to identifying or enforcing a complete material symmetry group. No completeness or universal-approximation theorem is asserted for a finite collection of these features followed by a monotone convex core.

\subsection{Reference normalization, growth, and polyconvexity}
Let $H(\bm z)$ denote a differentiable convex function that is nondecreasing in every feature, absorbing the positive scaling and centering into its definition. Set
\begin{equation}
 h_i=\partial_{z_i}H(\bm z^0)\geq0,\qquad r=\sum_i h_i\rho_i.
\end{equation}
The energy is
\begin{align}
 W(\F)={}&H(\bm z)-H(\bm z^0)-r(J-1)
 +\alpha(J-1-\log J)+\frac\beta2(J-1)^2\nonumber\\
 &+\varepsilon\left(\frac{\tr\C}{2}-1-\log J\right),
 \qquad \alpha,\beta,\varepsilon>0.
 \label{eq:energy}
\end{align}
Affine determinant stress normalization has an established precedent in Linden et al.~\cite{linden2023}. The role of Eq.~\eqref{eq:referencefeatures} here is to make that scalar normalization compatible with this anisotropic paired family. The sign of $r$ is unrestricted. Reference normalization is separated from the independent positive barrier coefficient $\alpha$ and quadratic coefficient $\beta$.

\begin{proposition}
Equation~\eqref{eq:energy} is objective, is polyconvex in two dimensions, and satisfies $W(\I)=0$ and $\partial_{\ee}W(\bm0)=\bm0$. It grows without bound as $J\to0^+$, as $J\to\infty$, and as $\|\F\|\to\infty$ at bounded positive $J$.
\end{proposition}
\begin{proof}
The composition $H(\bm z(\F,J))$ is convex in independent $(\F,J)$ because each $z_i$ is convex and $H$ is convex and nondecreasing. The correction $-r(J-1)$ is affine in $J$ for either sign of $r$. Every remaining nonconstant term is convex in $(\F,J)$. In two dimensions, convex dependence on $(\F,J)$ supplies the required convex representation in deformation minors. Objectivity follows from the norms and determinant. Equation~\eqref{eq:referencefeatures}, together with $\partial_{\ee}J|_0=(1,1,0)^T$, gives exact reference normalization. Finally, the supporting plane of $H$ and $z_i\geq0$ give a constant lower bound for $H$; the positive logarithmic, quadratic, and $\|\F\|^2$ terms dominate the affine correction in the stated limits.
\end{proof}

These are constitutive statements. They do not prove uniqueness of a structural equilibrium, absence of geometric buckling, or convergence of a particular Newton algorithm. Nor do they prove that a homogenized material with an arbitrary microscopic instability can be accurately represented by a globally polyconvex law.

\subsection{ICNN and integrated-hat spline cores}
The ICNN uses softplus activations and nonnegative hidden, input, and output weights. Convexity of ICNNs follows from the composition construction described by Amos et al.~\cite{amos2017}. Nonnegativity of the direct-input weights is an additional restriction needed here for monotonicity in the nonlinear features; convexity alone does not require every direct-input weight to be nonnegative.

For the spline core, each scalar edge has the form
\begin{equation}
 f(x)=a x+\sum_k c_k B_k(x)+d,\qquad a,c_k\geq0,
 \label{eq:splineedge}
\end{equation}
where $B_k''$ is a nonnegative triangular hat. With $s=(x-k)/h$,
\begin{equation}
 B_k(x)=h^2\begin{cases}
 0,&s\leq-1,\\
 (s+1)^3/6,&-1<s\leq0,\\
 (-s^3+3s^2+3s+1)/6,&0<s<1,\\
 s,&s\geq1.
 \end{cases}
 \label{eq:hat}
\end{equation}
Thus $f'$ and $f''$ are nonnegative, and $f''$ is continuous. Positive sums and layer compositions preserve convexity and monotonicity. Linear tails match analytically. Unlike a positive cubic-hinge construction, the curvature can decrease back to zero without becoming negative. This is a local $C^2$ spline implementation inspired by the monotonic input-convex KAN framework~\cite{thakolkaran2025}; it is not an unmodified reproduction of that paper's spline implementation or training procedure.

The distinction at the tails is mathematical rather than terminological. Appendix A.3 of Ref.~\cite{thakolkaran2025} specifies endpoint slopes using finite secants. For a strictly convex interior curve, a forward secant at the left endpoint is larger than the right derivative there, so that using it as a left-tail slope creates a downward derivative jump. An analogous issue occurs at the right endpoint. Exact tangent slopes would preserve convexity, but a finite-secant formula does not generally guarantee exact convexity or $C^1$ continuity across the joins. This observation concerns extrapolation beyond the natural spline interval; it does not negate the interior constrained-spline argument. Equation~\eqref{eq:hat} avoids this distinction through exact matching, with vanishing curvature at both tails.

\subsection{A separate nonnegative-energy check}
\label{sec:nonnegative}
Zero reference energy and stress do not, on their own, establish global nonnegative energy. This distinction is also explicit in Linden et al.~\cite{linden2023}. For Eq.~\eqref{eq:energy}, define $k_i=1/p_i+1/q_i$ and $\eta_i=\rho_i/k_i$. The determinant inequality $\|\F\bm d_i\|^2\|\F\bm d_i^\perp\|^2\geq J^2$ gives $z_i\geq k_iJ^{\eta_i}$. Together with $\tr\C/2\geq J$ and the supporting plane of $H$, this yields
\begin{equation}
 W\geq g(J):=\sum_i h_i k_i[J^{\eta_i}-1-\eta_i(J-1)]
 +(\alpha+\varepsilon)(J-1-\log J)+\frac\beta2(J-1)^2.
 \label{eq:lowerbound}
\end{equation}
Since $g(1)=g'(1)=0$, proving $g''\geq0$ for $J>0$ is sufficient. Appendix~\ref{app:nonnegative} describes the implemented bound. This is a saved-parameter certificate with numerical margins, not a restriction enforced at every training iterate and not a formally verified interval-arithmetic proof.

\section{Reduced microscopic models}
\label{sec:rom}
\subsection{Linear projection and fixed empirical cubature}
Let $\bm a\approx\bm V\qq$ approximate the independent periodic fluctuation coordinates in a POD basis. The linear reduced equilibrium is $\bm V^T\bm r(\bm V\qq,\ee)=0$. With elementwise contributions $\bm r_e$, an element-selected empirical cubature approximation has the form
\begin{equation}
 \bm r_r(\qq,\ee)=\sum_{e\in\mathcal Z_r}w_e\bm V_e^T
 \bm r_e(\bm V\qq,\ee)=\bm0,\qquad w_e\geq0.
 \label{eq:ecm}
\end{equation}
The selected contributions retain their underlying quadrature rules. The fixed ECM approach follows the reduced-integrand cubature principle of Hern\'andez et al.~\cite{hernandez2017}; the present implementation selects elements and deploys genuine reduced MDPA meshes. No neural network enters this HPROM tier.

Stress is computed with a separate fixed rule $\mathcal Z_s$, averaging microscopic first Piola stress before the macroscopic conversion in Eq.~\eqref{eq:homstress}. This separate output approximation is important: even when Eq.~\eqref{eq:ecm} derives from a fixed weighted microscopic energy, the separately integrated output is not automatically the derivative of that same reduced energy. We therefore do not assign an exact macroscopic potential guarantee to the deployed HPROM simply because it evaluates the original microscopic constitutive law.

\subsection{Nonlinear manifold and adaptive weights}
For the iterative HPROM--ANN, a decoder $\bm a=\bm g(\qq)$ replaces the linear subspace. A representative latent closure has the structure
\begin{equation}
 \bm g(\qq)=\bm V_M\qq+\bm V_S\bm n(\qq),\qquad
 \bm G(\qq)=\frac{\partial\bm g}{\partial\qq}.
\end{equation}
The microscopic residual is projected with the decoder tangent. Following the manifold-adaptive cubature setting~\cite{hernandez2026}, its approximation is
\begin{equation}
 \bm r_r(\qq,\ee)=\sum_{e\in\mathcal Z_r} w_e(\qq)\bm G_e(\qq)^T
 \bm r_e(\bm g(\qq),\ee)=\bm0.
 \label{eq:maw}
\end{equation}
The implementation uses positive softmax weight fields with a fixed prescribed sum. This guarantees the parametrized weight constraints at a query; it does not guarantee integration accuracy there. Stress has its own adaptive rule. Differentiation must include the decoder tangent, its variation where required, and the weight derivatives.

In particular, $\sum_e w_e(\qq)\partial_{\qq}\Pi_e$ is not generally equal to $\partial_{\qq}\sum_e w_e(\qq)\Pi_e$: the latter contains $\sum_e\Pi_e\partial_{\qq}w_e$. Structure-preservation results for fixed-weight ECSW and Lagrangian dynamical systems~\cite{farhat2015} must therefore not be transferred to Eq.~\eqref{eq:maw} without checking their hypotheses. A correct Jacobian is necessary for consistent Newton linearization but does not establish a common potential for independently approximated residuals and outputs.

\subsection{Direct closure}
The D-HPROM--ANN evaluates the slave amplitudes directly from the macroscopic engineering strain. In the deployed coordinates its displacement reconstruction has the form
\begin{equation}
 \bm a_{\rm direct}(\ee)=\bm V_{MA}\ee+\bm V_S\bm n(\ee).
\end{equation}
The microscopic equilibrium iterations are absent; the reconstructed state is used in the adaptive stress rule. The macroscopic tangent differentiates the complete direct stress map, including closure and weights, by central differences. This is a direct micromechanically evaluated surrogate, not an equilibrated reduced boundary-value problem. It retains local fields for stress evaluation but has no imposed macroscopic potential or polyconvexity guarantee.

\subsection{Removing online dependence on discarded elements}
The reduced implementations restrict lifting operators and decoder rows to the active MDPA degrees of freedom. Local force and stiffness contributions are projected directly, avoiding unnecessary full-dimensional assembly. Stress and tangent stencil evaluations are batched. For the direct model, decoder Jacobians that are unused by the finite-difference stress tangent are not computed. These are implementation changes to the same numerical models, not changes in supports, reduced dimensions, tolerances, or constitutive laws. Agreement with the retained original implementations and independent derivative checks is assessed before complete timing runs.

\section{Numerical example I: constitutive response of a periodic porous cell}
\label{sec:cell}
\subsection{Purpose and configuration}
This example asks whether a mechanically constrained learned law can represent the anisotropic response of the chosen cell, and which mechanical properties distinguish the learned alternatives. It does not measure structural speedups. The square periodic cell contains a centered elliptical void with area fraction $0.20$, aspect ratio $2$, and major-axis angle $30^\circ$ relative to the cell axes; see Fig.~\ref{fig:rve}. This orientation avoids the mirror-aligned configurations of the square cell. Centrosymmetry remains, but does not mean that the strain energy is an even function of strain.

The solid is compressible Neo-Hookean in plane strain:
\begin{equation}
 \psi_\mu=\frac\mu2(\tr\C_\mu-2)-\mu\log J_\mu
              +\frac\lambda2(\log J_\mu)^2,
\end{equation}
with Young's modulus $1.628$ GPa and Poisson ratio $0.4$. Here $\mu$ and $\lambda$ are the usual Lam\'e constants. This is the specific matrix law used by the reference solver, not a generic assertion about every Neo-Hookean variant. The production mesh contains 1546 quadratic triangles, 3240 nodes, and 4638 quadrature points. Periodicity and null-space constraints leave 6320 independent microscopic displacement coordinates. Cell averages use the complete reference cell measure.

\begin{figure}[htbp]
 \centering\includegraphics[width=.96\linewidth]{rve_geometry.pdf}
 \caption{Microstructure and production discretization. The left panel shows corner-edge traces of the quadratic mesh; the right panel specifies the void geometry. The cell side is denoted by $\ell$. No material reflection symmetry is imposed on the learned paired features.}
 \label{fig:rve}
\end{figure}

The production mesh was selected using a local stress-field criterion along an investigated uniaxial path, rather than only a homogenized average. At $E_{11}=0.20$, its maximum microscopic first Piola stress norm differs by approximately $0.486\%$ from a 3095-element reference, whereas the average axial stress differs by approximately $0.001\%$. This is a mesh check at a specified state, not a uniform error bound over all multiaxial data.

\subsection{Data domain and model selection}
The fitting domain is informed by a preliminary macroscopic deformation study. The final training box, shared by the constitutive and reduced models, has the following rounded bounds:
\begin{equation}
 \ee_{\min}=(0,-0.09865,-0.16807)^T,\qquad
 \ee_{\max}=(0.1955,0.01635,0.10793)^T.
 \label{eq:box}
\end{equation}
The new constitutive features do not change this box. Of 4950 available training states, 4208 are used for fitting and 742 for validation, with split seed 5. The independent test set contains 400 states. An additional 350-state probe extends selected distances from the box center by factors $1.05$, $1.10$, $1.25$, $1.50$, and $2.00$. Five probe labels are nonfinite following failed reference solves and are excluded with the same finite-label mask for all compared models. Their exclusion is reported, not interpreted as proof of a physical instability.

The selected feature-learning runs use full-batch Adam for 2600 steps followed by up to 400 L-BFGS iterations. The normalized training objective combines stress, energy, and reference-tangent relative mean-square errors with weights $1$, $0.2$, and $0.02$, respectively. The tangent target is obtained from a separate coarse periodic reference calculation. Candidates are selected by validation stress error; test and probe labels are not used in that second-stage selection. Both selected runs use initialization seed 16, with 32 features and hidden widths $24/24$ for ICNN and $12/12$ for ICKAN. Earlier baselines have different feature families, objectives, and training budgets. Comparisons to them diagnose practical approximability, not compute-matched superiority of one network architecture.

\subsection{From a restrictive feature family to an accurate fit}
The earlier 15-feature construction combines $\tr\C$, twelve balanced directional fourth- and sixth-power features of $\F$ and $\cof\F$, and $J,J^2$. Its reference correction uses $-r\log J$. The revised formulation changes both the feature family and the separation of reference normalization from volumetric stiffness; it is not simply a wider network or a relabelling of the old inputs.

A convex nondecreasing core preserves the componentwise ordering of its inputs. Consequently, if all features increase along a deformation path, the core output cannot decrease. Whether the complete energy decreases also depends on its analytic correction terms. An incompatibility with a fixed feature/correction family is a representational restriction of that family, not a contradiction of polyconvexity as a mathematical framework. The revised directional powers and determinant exponents provide additional ways to capture the observed deformation response while maintaining the sufficient convexity conditions.

The saved original ICNN and ICKAN have held-out aggregate stress errors of approximately $22.78\%$ and $18.33\%$, respectively. The selected revised models reduce these errors to approximately $0.65\%$; Table~\ref{tab:constitutive} reports the independently audited selected checkpoints. The Free model's corresponding aggregate stress errors are $1.2643\%$ on the test set and $5.5024\%$ on the finite probe set. Thus the polyconvex models improve this test fit but do not outperform Free on aggregate probe error. This distinction matters when discussing extrapolation.

\begin{table}[htbp]
 \centering\small
 \caption{Selected polyconvex models: constitutive errors, in percent.}
 \label{tab:constitutive}
 \begin{tabular}{lrrr}
\toprule
Model & $e_S^{\rm test}$ & $e_W^{\rm test}$ & Worst test \\
\midrule
ICNN & 0.6491 & 0.2169 & 5.53 \\
ICKAN & 0.6551 & 0.2148 & 5.78 \\
\bottomrule
\end{tabular}

\par\smallskip{\footnotesize All values are percentages. Stress and energy errors are aggregate array norms; the last column is the maximum per-state stress error. Material A: 400 independent test states.}

\end{table}

\begin{figure}[htbp]
 \centering\includegraphics[width=.85\linewidth]{probe_errors.pdf}
 \caption{Stress error on increasingly distant probe rings, using only finite reference labels. Ring factors expand the sampling box; they are not coupon load multipliers. The increase to approximately $12.25\%$ at factor 2 illustrates a remaining extrapolation limitation despite architectural admissibility.}
 \label{fig:probe}
\end{figure}

The archived independent audit records successful energy first- and second-derivative checks and saved-weight nonnegative-energy sufficient bounds for both selected models. The conservative finite-interval margins are $17.25$ and $16.82$ in normalized energy units. Random admissible-deformation tests support implementation consistency, but the polyconvexity argument is analytical. Neither aggregate error nor a successful bound on energy is a uniform bound on stress prediction error.

\subsection{A closed-cycle witness for direct stress regression}
\label{sec:cycle}
For a differentiable hyperelastic potential, any closed strain cycle in its domain satisfies
\begin{equation}
 \mathcal W_{\rm cyc}=\oint\ssv\cdot\dd\ee=0.
 \label{eq:cycle}
\end{equation}
This test isolates integrability and does not require a structural Newton failure. The importance of conservative strain--stress mappings is discussed explicitly by As'ad et al.~\cite{asad2022}. We evaluate a counterclockwise rectangular cycle in $(E_{11},E_{22})$, with half-width $0.001$, centered at
\begin{equation}
 \ee_c=(0.0051260033,-0.0945507613,-0.1283720602)^T.
\end{equation}
The engineering shear component is held fixed. The rectangle lies inside Eq.~\eqref{eq:box}. Its center was chosen after inspecting regression tangent asymmetry at held-out states; it is a deliberately selected mechanical witness, not an unbiased estimate of the frequency of violations. The models remain frozen during this audit.

\begin{figure}[htbp]
 \centering\includegraphics[width=.97\linewidth]{cycle_convergence.pdf}
 \caption{Closed-cycle geometry and quadrature refinement of net stress work. Regression approaches a nonzero signed value of $-213.922$ J m$^{-3}$. Free and ICNN cancel to numerical precision; an exact numerical zero is plotted at $10^{-12}$ J m$^{-3}$ solely for logarithmic display. The $C^2$ spline energy also defines a conservative law, but crossing spline knots requires quadrature refinement; its residual work decreases towards zero. The FOM marker is a separately evaluated cycle at eight Gauss points per edge.}
 \label{fig:cycle}
\end{figure}

Table~\ref{tab:cycle} retains the signed work and its quadrature dependence. Reversing the cycle reverses the regression work, which is incompatible with interpreting that map as an elastic potential. The result establishes a constitutive inconsistency, not by itself a quantitatively important error for every macroscopic loading history. The Free model passes this test because it is energy-based; energy-cycle consistency therefore cannot distinguish Free from the polyconvex models.
\begin{table}[htbp]
 \centering\small
 \caption{Closed-cycle work under quadrature refinement.}
 \label{tab:cycle}
 \begin{tabular}{lrrrrr}
\toprule
Model & 8 & 16 & 32 & 64 & 128 \\
\midrule
Regression & -213.922 & -213.922 & -213.922 & -213.922 & -213.922 \\
Free & -5.82077e-11 & 1.05501e-10 & 1.33241e-10 & -7.83302e-11 & 3.57659e-10 \\
ICNN & 6.54836e-11 & -3.81988e-11 & -1.30967e-10 & -1.59162e-12 & 4.0157e-10 \\
ICKAN & 0.117515 & -0.0171476 & -0.00114144 & -0.000254375 & 1.08569e-05 \\
HPROM & -1.23416 & -1.23416 & -1.23416 & -1.23416 & -1.23416 \\
HPROM--ANN & 247.774 & 247.774 & 247.774 & 247.774 & 247.774 \\
D-HPROM--ANN & 54.0815 & 54.0815 & 54.0815 & 54.0815 & 54.0815 \\
\bottomrule
\end{tabular}

\par\smallskip{\footnotesize Signed cycle work in J m$^{-3}$; column headings give Gauss points per edge. FOM at 8 points per edge: -6.98492e-10 J m$^{-3}$. The spline quadrature error decreases with refinement; it is not physical dissipation. Reversal and equilibrium-tolerance controls are discussed in the text.}

\end{table}

\subsection{Rank-one curvature: guarantees, samples, and an unresolved reference}
For $\bm H=\bm a\otimes\bm b$, the relevant finite-strain curvature is
\begin{equation}
 D^2_{\F}W[\bm H,\bm H]
   =(\delta\ee)^T\bm D\,\delta\ee+\Sg:(\bm H^T\bm H),\quad
 \delta\E=\operatorname{sym}(\F^T\bm H),
 \label{eq:lh}
\end{equation}
where $\delta\ee$ uses Eq.~\eqref{eq:voigt}. Omitting the second term tests a different property. In particular, a positive semidefinite tangent with respect to $\ee$ does not generally imply rank-one convexity in compression. For energy-based models, Eq.~\eqref{eq:lh} is evaluated over sampled states and unit directions; negative candidates are checked by second differences of $W(\F\pm h\bm H)$.

\begin{table}[htbp]
 \centering\small
 \caption{Finite rank-one curvature audit. Positive sampled values do not establish a global property.}
 \label{tab:rankone}
 \begin{tabular}{lrrrr}
\toprule
State set & $N$ & Free & ICNN & ICKAN \\
\midrule
Held-out test & 400 & 176.681 & 220.486 & 220.468 \\
Finite-label probe & 345 & 113.034 & 203.280 & 203.538 \\
In-box audit (unlabelled) & 12000 & 163.040 & 213.979 & 211.168 \\
\bottomrule
\end{tabular}

\par\smallskip{\footnotesize Minimum sampled rank-one curvature in MPa. Minima need not occur at the same state or direction. The additional in-box cloud has no independently computed FOM labels; 180 unit directions for $\bm b$ are sampled, with minimization over $\bm a$ via an eigenproblem. None of these finite samples is a global certificate.}

\end{table}

No negative direction is found for Free in the audited test, finite-label probe, and additional in-box sets. An outward search over 1800 strain-space rays finds its first sampled negative value at expansion factor $1.55$, with $J\simeq0.780$ and curvature approximately $-3.43$ MPa. This is a result of that search, not the exact nearest violation. Independent energy finite differences confirm the negative curvature of the learned model.

However, the attempted FOM continuation fails numerically before reaching that candidate. No refined-continuation or global branch-tracking result is available there. Therefore this candidate establishes that the Free network lacks global rank-one convexity, but does not establish disagreement with a converged stable homogenized reference at that state. Solver nonconvergence is not evidence that the physical branch ceases to exist. The stronger claim of an FOM-confirmed false instability is left open. The ICNN and ICKAN have an analytical polyconvexity argument; their sampling audit is a check of implementation, not its replacement.

\section{Numerical example II: structural accuracy and online cost}
\label{sec:coupon}
\subsection{Purpose, boundary conditions, and coverage}
This example asks how the constitutive approximations affect a complete nonlinear structural calculation. The material is the same periodic cell as in Section~\ref{sec:cell}; the new evidence concerns error propagation and online cost, not a second microstructure. The dog-bone profile has overall length 165 mm, grip width 19 mm, narrow-section width 13 mm, narrow-section length 57 mm, and fillet radius 76 mm. It is a plane-strain computational benchmark using an ASTM-inspired in-plane profile, not a simulation certified to reproduce an ASTM tensile test.

The two end faces receive equal-and-opposite axial resultants of 100 kN per end, applied in 20 equal increments. Three point constraints remove rigid modes: the axial and transverse displacement at the left mid-edge node, and transverse displacement at the right mid-edge node. Entire end faces are not clamped. Consequently, a small skew of the deformed ends is not ruled out by the boundary conditions. The anisotropic response can couple axial extension to shear.

The macro mesh contains 690 quadratic triangles, 1501 nodes, and 2070 integration points. The reference extrusion thickness used to convert the two-dimensional stress integrals into the stated force resultants is 50 mm. This force convention is part of the plane-strain benchmark, not an ASTM specimen-thickness specification. The physical cell size is declared as 1.3 mm, corresponding to ten cells across the narrow width. This provides a scale-separation assumption, not a direct numerical validation of homogenization error at that ratio. All saved converged integration-point strains for every reported model remain inside the unchanged training box. This check does not automatically cover unsaved intermediate Newton iterates or certify another load case.

\begin{figure}[htbp]
 \centering\includegraphics[width=.99\linewidth]{coupon_fields.pdf}
 \caption{Nested FOM--FE$^2$ reference at 100 kN per end. Fields are displayed on the actual deformed mesh, without displacement amplification. Gauss-point strain and stress values are first averaged within each element and then patch-averaged to nodes for visualization only. These recovered contours are not the raw fields used in error norms and can smooth local extrema. Stresses are second Piola components. A three-dimensional von Mises stress is not reconstructed because the saved output does not include the necessary out-of-plane Cauchy stress.}
 \label{fig:coupon}
\end{figure}

\subsection{Online models and implementation verification}
The linear HPROM uses 39 reduced coordinates, a residual rule supported on 135 elements, and a stress rule on 73 elements; their union contains 183 elements. The iterative HPROM--ANN uses three latent coordinates and 36 learned slave amplitudes, with ten residual and ten stress elements and a union of 19. The direct model evaluates stress on ten elements and solves no microscopic equilibrium equations. These counts refer to actual reduced MDPA meshes, not masked evaluations on the full RVE mesh. Figure~\ref{fig:supports} shows their spatial distribution and the overlap of residual and output rules.

\begin{figure}[htbp]
 \centering\includegraphics[width=.99\linewidth]{reduced_supports.pdf}
 \caption{Element supports of the deployed hyperreduced microscopic models. The pale full-cell mesh provides location context only; discarded elements are not assembled online. Colors distinguish residual-only, stress-only, and shared elements. Supports are read from the same index arrays used to construct the reduced MDPA meshes.}
 \label{fig:supports}
\end{figure}

For the optimized linear and direct implementations, the maximum normalized stress differences from the original implementations over the verification states are $4.71\times10^{-14}$ and $4.20\times10^{-13}$. Corresponding tangent differences are $2.68\times10^{-10}$ and $7.25\times10^{-10}$. Independent finite-difference checks give discrepancies of approximately $3.51\times10^{-9}$ and $2.38\times10^{-9}$. Such comparisons check that optimization preserves the deployed model; they do not establish that either model exactly reproduces full microscopic equilibrium.

\subsection{Error and timing definitions}
For a saved nodal or integration-point array $A$, the reported relative error is
\begin{equation}
 e_A=100\frac{\|\operatorname{vec}(A-A_0)\|_2}{\|\operatorname{vec}(A_0)\|_2},
 \label{eq:error}
\end{equation}
where $A_0$ is the nested reference on the same meshes. These are unweighted discrete array norms. The strain array stores engineering shear, whereas the stress array stores tensor shear. They are neither quadrature-weighted continuum $L^2$ norms nor invariant tensor Frobenius norms. The metric is kept identical across models; this does not make the three different field errors interchangeable.

Online wall time includes the complete macroscopic solve, microscopic or surrogate stress/tangent evaluations, communication, and lazy worker construction within the timed solve. Parent setup, offline data generation, training, reduced-mesh construction outside the timed region, and result export are excluded. The FOM uses 20 worker processes with one numerical-library thread per worker. The three optimized microscopic surrogates use the same process configuration and three sequential fresh complete runs. PANN evaluations are batched in one process and use the separately investigated thread settings shown in Table~\ref{tab:fe2}.

The macroscopic solve stops when the free-degree-of-freedom residual norm falls below $10^{-5}$ N or its ratio to the initial increment residual falls below $10^{-7}$. The periodic FOM uses a relative microscopic displacement-increment tolerance of $10^{-10}$, with continuation when needed. The archived FOM and selected PANN times each have one repetition in this comparison. Thus their ratios are observed online speedups under the stated implementations, not matched-thread efficiency estimates or repeated-run confidence intervals. Table~\ref{tab:repeats} exposes the available reduced-model variability rather than selecting its fastest repetition. A publication-final timing study should repeat every tier under a common resource-budget protocol and retain all repetitions.

\begin{table}[htbp]
 \centering\footnotesize
 \caption{Complete coupon simulations: final-field errors against the nested FOM and archived online wall times.}
 \label{tab:fe2}
 \begin{tabular}{lrrrrrrr}
\toprule
Model & CPU & $n_{\rm run}$ & $t$ [s] & $t_0/t$ & $e_u$ [\%] & $e_E$ [\%] & $e_S$ [\%] \\
\midrule
FOM & 20p & 1 & 4268.679 & 1.0 & 0.00000 & 0.00000 & 0.00000 \\
HPROM & 20p & 3 & 67.348 & 63.4 & 0.00565 & 0.01653 & 0.00099 \\
HPROM--ANN & 20p & 3 & 30.143 & 141.6 & 0.05200 & 0.14968 & 0.03727 \\
D-HPROM--ANN & 20p & 3 & 3.788 & 1126.9 & 0.02324 & 0.04851 & 0.00967 \\
ICNN & 4t & 1 & 1.734 & 2461.3 & 0.42468 & 0.77028 & 0.08880 \\
ICKAN & 8t & 1 & 3.669 & 1163.5 & 0.38657 & 0.78932 & 0.07800 \\
Free & 12t & 1 & 1.356 & 3146.9 & 0.92988 & 1.49464 & 0.13492 \\
Regression & 16t & 1 & 0.946 & 4510.4 & 0.60620 & 0.81218 & 0.07378 \\
\bottomrule
\end{tabular}

\par\smallskip{\footnotesize p: worker processes, one numerical-library thread each; t: PyTorch threads. For $n_{\rm run}=3$, time is the median; for $n_{\rm run}=1$, it is the archived single run. These are online wall times, not offline-inclusive or matched-thread speedups.}

\end{table}
\begin{table}[htbp]
 \centering\small
 \caption{Sequential optimized-run repetitions used for median HPROM timings.}
 \label{tab:repeats}
 \begin{tabular}{lrrr}
\toprule
Model & $t_1$ [s] & $t_2$ [s] & $t_3$ [s] \\
\midrule
HPROM & 63.354 & 67.348 & 68.953 \\
HPROM--ANN & 27.947 & 34.628 & 30.143 \\
D-HPROM--ANN & 3.788 & 3.758 & 3.969 \\
\bottomrule
\end{tabular}

\end{table}

\subsection{Accuracy--cost trade-off and Newton behavior}
Figure~\ref{fig:cost} and Table~\ref{tab:fe2} show distinct operating points rather than a single universally best model. The linear HPROM gives the smallest displacement error, approximately $0.00565\%$, at a median online time of 67.35 s. The iterative HPROM--ANN reduces this time to 30.14 s with a displacement error of approximately $0.0520\%$. The optimized iterative nonlinear model is therefore about $2.23$ times faster than the \emph{optimized} linear model; comparison against the old linear implementation would overstate this gain. The direct closure reduces the time to 3.79 s and has a displacement error of approximately $0.0232\%$ on this path. Its lower error than the iterative nonlinear model is an observed outcome, not a consequence of a general error ordering.

The ICNN and ICKAN times are 1.73 s and 3.67 s, with displacement errors of approximately $0.425\%$ and $0.387\%$. Their finite-strain structural accuracy is therefore consistent with a useful constitutive approximation, although it is poorer than the reduced microscopic models on this benchmark. Direct stress regression is fastest, but its successful coupon solution does not remove the integrability defect exposed by Section~\ref{sec:cycle}. Free also converges and remains a relevant baseline; its absence of a global convexity constraint is not equivalent to a failure on this load case.

\begin{figure}[htbp]
 \centering\includegraphics[width=.93\linewidth]{cost_accuracy.pdf}
 \caption{Online cost versus final displacement error on the common coupon problem. Horizontal segments show the observed range of three optimized HPROM runs; unbarred points are archived single PANN runs. Process and thread counts differ as detailed in Table~\ref{tab:fe2}. The plot does not include offline costs.}
 \label{fig:cost}
\end{figure}

All models require four macroscopic iterations per increment in the archived runs. Consequently, this example provides no empirical evidence that polyconvexity reduced Newton iterations relative to Free or Regression. It establishes that the selected constrained models can be used successfully with the common structural solver, while maintaining their constitutive guarantees. A separate challenge problem would be needed to establish a solver-robustness advantage, and its reference solution would need to be validated independently.

\section{Position relative to prior work and limitations}
\label{sec:discussion}
\subsection{A formulation-specific literature comparison}
Table~\ref{tab:literature} states what the selected sources actually formulate, distinguishing variants within a paper. It does not use a missing demonstration as evidence that a method is incapable of a property. Page and equation locations used to check these statements are retained in the accompanying source audit. The full supplied bibliography remains under review in this working draft; the table is not an exhaustive systematic review.

\begin{table}[htbp]
\centering\footnotesize
\caption{Scope of selected antecedents. ``By construction'' refers to the stated formulation and hypotheses, not all implementations sharing an architecture name.}
\label{tab:literature}
\begin{tabularx}{\linewidth}{>{\raggedright\arraybackslash}p{2.2cm}YYY}
\toprule
Source / variant & Constitutive or reduced construction & Property relevant here & Distinction retained in this paper\\
\midrule
As'ad et al.~\cite{asad2022} & Energy learned from stress derivatives; ICNN in Green strain & Objective, conservative energy; convexity in strain & Tensile-state stability argument is not a global polyconvexity theorem\\
Klein et al.~\cite{klein2022}, invariant model & Convex monotone network of polyconvex invariants & Objectivity and polyconvexity by construction & Specific invariant choices can restrict approximation\\
Klein et al.~\cite{klein2022}, gradient model & Network of deformation minors and group symmetrization & Polyconvexity; material group enforced & Objectivity approximated by data augmentation\\
Linden et al.~\cite{linden2023} & Invariant PANN with compatible analytic corrections & Polyconvex reference normalization and growth for stated anisotropies & Global nonnegative energy is treated separately\\
Klein et al.~\cite{klein2026}, triclinic model & Objective polyconvex features of $\C$, $\cof\C$, and $\pm J$ & Optional finite-group symmetrization; can omit known material symmetry & Not restricted in formulation to its cubic numerical examples\\
Thakolkaran et al.~\cite{thakolkaran2025} & Monotonic input-convex spline KAN & Polyconvex isotropic construction on admissible spline domains & Finite-secant tail joins need qualification; our integrated-hat core matches tails analytically\\
Farhat et al.~\cite{farhat2015} & Fixed-weight ECSW for reduced nonlinear dynamics & Lagrangian structure preservation under its hypotheses & Does not automatically cover our adaptive weights and independent stress rule\\
Hern\'andez et al.~\cite{hernandez2026} & Empirical cubature with manifold-adaptive weights & Reduced nonlinear equilibrium with adaptive integration & Existing cubature methodology used here, not a new constitutive certificate\\
\bottomrule
\end{tabularx}
\end{table}

\subsection{What the present evidence supports}
The principal positive result is practical compatibility between accurate fitting and a polyconvex architecture for the investigated anisotropic cell. The original training difficulty does not show that polyconvexity is unreliable; it shows that sufficient constructions differ in expressive restrictions. The paired features and independent volume terms substantially improve the fit, but the experiment changes several ingredients simultaneously. A frozen-budget ablation is still needed to attribute the gain separately to feature learning, determinant exponents, reference normalization, and spline curvature flexibility.

The mechanical tests address different requirements. Closed-cycle work detects the lack of a potential in Regression. Rank-one curvature addresses an additional condition not implied by the potential structure of Free. The current Free counterexample is outside the reference-validated region and is not an FOM-confirmed artificial instability. Nominal coupon convergence therefore complements these tests rather than replacing them. It would be scientifically incorrect to report that all unconstrained alternatives fail while all constrained alternatives are accurate everywhere.

The results concern a single material, plane-strain kinematics, one structural loading path, and a prescribed training box. They do not demonstrate universality across microstructures, transfer of one trained network to another material, three-dimensional constitutive completeness, or history-dependent behavior. A second microstructure should answer an additional question---for example, whether the feature family remains effective under qualitatively different deformation mechanisms---rather than simply duplicate the same tensile test. Such a result is not included before its data exist.

Finally, the current online comparison excludes potentially substantial offline costs and has unequal repetition counts and thread configurations. These limitations are visible in the tables. Break-even query counts, deployment robustness under domain exit, and a common repeated timing protocol are necessary complements for a publication-final assessment of computational efficiency.

\section{Conclusions}
Paired directional features with constrained trainable powers and determinant exponents provide a practical polyconvex constitutive construction for the anisotropic periodic material studied here. Their isotropic reference derivatives allow exact stress normalization through an affine determinant correction without tying the normalization coefficient to the volumetric barrier stiffness. Both a smooth ICNN and an integrated-hat spline-ICKAN achieve approximately $0.65\%$ held-out aggregate stress error and below $0.43\%$ displacement error in the nominal coupon simulation.

Optimized reduced microscopic models occupy a different accuracy--cost range. They are more accurate on this coupon, while direct learned constitutive evaluation is generally faster. A closed-cycle test exposes a conservative-mechanics limitation of the regression baseline that nominal finite element convergence does not reveal. In contrast, the available Free-model curvature witness is not yet supported by a converged FOM reference and is reported only within that limited scope. The study supports the usefulness of the proposed constrained feature family for the present problem, not unrestricted extrapolation or a universal ranking of constitutive and reduced-order models.

\appendix
\section{Joint convexity of the directional building block}
\label{app:convex}
Consider $f(\bm x,J)=\|\bm x\|^aJ^{-b}$ with $J>0$, $a\geq1$ and $0\leq b\leq a-1$. For $t=\|\bm x\|>0$, the Hessian block in the radial and determinant variables is
\begin{equation}
 \begin{pmatrix}
 a(a-1)t^{a-2}J^{-b} & -ab t^{a-1}J^{-b-1}\\
 -ab t^{a-1}J^{-b-1} & b(b+1)t^aJ^{-b-2}
 \end{pmatrix}.
\end{equation}
Its diagonal entries are nonnegative and its determinant is
\begin{equation}
 ab(a-b-1)t^{2a-2}J^{-2b-2}\geq0.
\end{equation}
The transverse Hessian eigenvalues are $a t^{a-2}J^{-b}\geq0$. Continuous extension gives convexity at $\bm x=0$. For the constitutive evaluation on $\GL$, a nonzero reference direction cannot satisfy $\F\bm d=0$, so this point does not introduce an interior singularity of the directional norm. Taking $a=2p$ or $2q$ establishes Eq.~\eqref{eq:featureconstraints}.

\section{Sufficient lower bound for nonnegative energy}
\label{app:nonnegative}
Write $t=\|\F\bm d_i\|^2$, $u=\|\F\bm d_i^\perp\|^2$. Hadamard's determinant inequality gives $tu\geq J^2$. Since each feature is increasing in $t,u$, its minimum at fixed $J$ occurs on $tu=J^2$. Minimizing
\begin{equation}
 \frac{t^{p_i}}{p_iJ^{b_i}}+\frac{J^{2q_i-c_i}t^{-q_i}}{q_i}
\end{equation}
gives $t^{p_i+q_i}=J^{2q_i+b_i-c_i}$ and hence $z_i\geq k_iJ^{\eta_i}$. This is a separate lower bound for each feature; simultaneous attainability for all directions is not assumed. It therefore remains valid after summing the supporting-plane inequalities.

For Eq.~\eqref{eq:lowerbound},
\begin{equation}
 J^2g''(J)=\alpha+\varepsilon+\beta J^2+\sum_i c_i^*J^{\eta_i},
 \qquad c_i^*=h_i k_i\eta_i(\eta_i-1).
\end{equation}
Only terms with $0<\eta_i<1$ can have negative coefficients. If $N=\sum_{c_i^*<0}|c_i^*|$, the two inequalities
\begin{align}
 \alpha+\varepsilon+\sum_{\eta_i<0}c_i^*-N&\geq0,\\
 \beta+\sum_{\eta_i\geq1}c_i^*-N&\geq0
\end{align}
are sufficient on $0<J\leq1$ and $J\geq1$, respectively. Their failure is inconclusive. The saved-parameter audit additionally uses lower bounds on powers over finite determinant intervals and analytic bounds on both infinite tails. Positive margins establish this sufficient analytical argument for the saved parameters to the accuracy of its floating-point evaluation; random energy samples alone would not establish it.

\section{Reproducibility and draft status}
The accompanying \texttt{build\_evidence.py} reads archived result arrays and metadata, verifies common mesh hashes, loads, completed increments, and saved-state coverage, and regenerates the numerical tables and vector figures. The generated \texttt{evidence\_manifest.json} records SHA-256 hashes of the consumed artifacts. It performs no retraining or new timing run. Selected checkpoint hashes and training manifests are retained in the constitutive audit directory. Thus regenerated graphics cannot silently be interpreted as newly measured performance.

The source text of the earlier manuscript is preserved. This draft changes its organization and narrows claims to available evidence. Its accompanying \texttt{REFERENCE\_AUDIT.md} distinguishes extraction, inspected source passages, and unfinished full-paper reading. The supplied Boehler chapter is incomplete, and no theorem is attributed to its missing pages. Author metadata, complete reference review, final statistical timing protocol, and any additional numerical example remain editorial or experimental completion tasks rather than implicit completed results.

\begin{thebibliography}{99}
\bibitem{yvonnet2007} J. Yvonnet, Q.-C. He, The reduced model multiscale method (R3M) for the non-linear homogenization of hyperelastic media at finite strains, Journal of Computational Physics 223 (2007) 341--368. \doi{10.1016/j.jcp.2006.09.019}
\bibitem{hernandez2017} J.A. Hern\'andez, M.A. Caicedo, A. Ferrer, Dimensional hyper-reduction of nonlinear finite element models via empirical cubature, Computer Methods in Applied Mechanics and Engineering 313 (2017) 687--722. \doi{10.1016/j.cma.2016.10.022}
\bibitem{barnett2023} J. Barnett, C. Farhat, Y. Maday, Neural-network-augmented projection-based model order reduction for mitigating the Kolmogorov barrier to reducibility, Journal of Computational Physics 492 (2023) 112420.
\bibitem{hernandez2026} J.A. Hern\'andez, S. Ares de Parga, R. Rossi, Dimensional hyperreduction of nonlinear finite element models via empirical cubature with manifold-adaptive weights, supplied manuscript (2026). Version consulted: local manuscript in \texttt{maw-ecm-paper}, 6 September 2026; publication status to be confirmed.
\bibitem{asad2022} F. As'ad, P. Avery, C. Farhat, A mechanics-informed artificial neural network approach in data-driven constitutive modeling, International Journal for Numerical Methods in Engineering 123 (2022) 2738--2759. \doi{10.1002/nme.6957}
\bibitem{kalina2023} K.A. Kalina, L. Linden, J. Brummund, M. K\"astner, FEANN: an efficient data-driven multiscale approach based on physics-constrained neural networks and automated data mining, Computational Mechanics 71 (2023) 827--851. \doi{10.1007/s00466-022-02260-0}
\bibitem{ball1977} J.M. Ball, Convexity conditions and existence theorems in nonlinear elasticity, Archive for Rational Mechanics and Analysis 63 (1977) 337--403.
\bibitem{klein2022} D.K. Klein, M. Fern\'andez, R.J. Martin, P. Neff, O. Weeger, Polyconvex anisotropic hyperelasticity with neural networks, Journal of the Mechanics and Physics of Solids 159 (2022) 104703. \doi{10.1016/j.jmps.2021.104703}
\bibitem{linden2023} L. Linden, D.K. Klein, K.A. Kalina, J. Brummund, O. Weeger, M. K\"astner, Neural networks meet hyperelasticity: A guide to enforcing physics, arXiv:2302.02403v2 (2023). \url{https://arxiv.org/abs/2302.02403}
\bibitem{klein2026} D.K. Klein, K.A. Kalina, R. Ortigosa, J. Mart\'inez-Frutos, M. K\"astner, O. Weeger, Advances in polyconvex anisotropic hyperelasticity, arXiv:2605.27011v1 (2026). \url{https://arxiv.org/abs/2605.27011}
\bibitem{thakolkaran2025} P. Thakolkaran, Y. Guo, S. Saini, M. Peirlinck, B. Alheit, S. Kumar, Can KAN CANs? Input-convex Kolmogorov-Arnold Networks (KANs) as hyperelastic constitutive artificial neural networks (CANs), Computer Methods in Applied Mechanics and Engineering 443 (2025) 118089. \doi{10.1016/j.cma.2025.118089}
\bibitem{amos2017} B. Amos, L. Xu, J.Z. Kolter, Input convex neural networks, Proceedings of Machine Learning Research 70 (2017) 146--155. \url{https://proceedings.mlr.press/v70/amos17b.html}
\bibitem{farhat2015} C. Farhat, T. Chapman, P. Avery, Structure-preserving, stability, and accuracy properties of the energy-conserving sampling and weighting method for the hyper reduction of nonlinear finite element dynamic models, International Journal for Numerical Methods in Engineering 102 (2015) 1077--1110. \doi{10.1002/nme.4820}
\end{thebibliography}
\end{document}
